\documentclass{article}
\usepackage{amsmath,amssymb,amsthm}
\usepackage{algorithm}
\usepackage[noend]{algpseudocode}
\usepackage{hyperref}
\usepackage{geometry}
\geometry{margin=1in}
\newtheorem{theorem}{Theorem}
\newtheorem{lemma}{Lemma}
\newtheorem{assumption}{Assumption}
\newcommand{\R}{\mathbb{R}}
\newcommand{\gap}{\mathcal{G}}
\newcommand{\prox}{\operatorname{prox}}

\begin{document}

\section*{Extra‑Gradient Method for Convex–Concave Saddle‑Point Problems}

\section*{Mathematical Formulation}
Let $L:\mathcal{X}\times\mathcal{Y}\rightarrow\R$ be a differentiable function that is convex in $x\in\mathcal{X}\subset\R^{n}$ and concave in $y\in\mathcal{Y}\subset\R^{m}$.  
We aim to solve the saddle‑point problem
\[
\min_{x\in\mathcal{X}}\;\max_{y\in\mathcal{Y}} L(x,y).
\]

Define the monotone operator
\[
F(z)\;:=\;\begin{pmatrix}
\nabla_{x}L(x,y)\\[2mm]
-\nabla_{y}L(x,y)
\end{pmatrix},\qquad 
z:=\begin{pmatrix}x\\ y\end{pmatrix}\in\mathcal{Z}:=\mathcal{X}\times\mathcal{Y}.
\]

Given a stepsize $\eta>0$, the \textbf{Extra‑Gradient (EG)} iteration reads
\begin{align}
z_{k+\frac12}&=z_{k}-\eta\,F(z_{k})\label{eq:mid}\\
z_{k+1}&=z_{k}-\eta\,F\!\bigl(z_{k+\frac12}\bigr).\label{eq:update}
\end{align}
In coordinates,
\[
\begin{aligned}
x_{k+\frac12}&=x_{k}-\eta\,\nabla_{x}L(x_{k},y_{k}), &
y_{k+\frac12}&=y_{k}+\eta\,\nabla_{y}L(x_{k},y_{k}),\\[1mm]
x_{k+1}&=x_{k}-\eta\,\nabla_{x}L(x_{k+\frac12},y_{k+\frac12}), &
y_{k+1}&=y_{k}+\eta\,\nabla_{y}L(x_{k+\frac12},y_{k+\frac12}).
\end{aligned}
\]

\section*{Pseudocode}
\begin{algorithm}[H]
\caption{Extra‑Gradient Method}
\begin{algorithmic}[1]
\Require Initial point $z_{0}=(x_{0},y_{0})$, stepsize $\eta>0$, max iterations $K$.
\For{$k=0,1,\dots,K-1$}
    \State Compute the gradient at the current point:
    \[
    g_{k}=F(z_{k})=
    \begin{pmatrix}
    \nabla_{x}L(x_{k},y_{k})\\ -\nabla_{y}L(x_{k},y_{k})
    \end{pmatrix}.
    \]
    \State \textbf{Predictor step} (mid‑point):
    \[
    z_{k+\frac12}=z_{k}-\eta\,g_{k}.
    \]
    \State Compute the gradient at the predictor:
    \[
    \tilde g_{k}=F(z_{k+\frac12})=
    \begin{pmatrix}
    \nabla_{x}L(x_{k+\frac12},y_{k+\frac12})\\ -\nabla_{y}L(x_{k+\frac12},y_{k+\frac12})
    \end{pmatrix}.
    \]
    \State \textbf{Corrector step}:
    \[
    z_{k+1}=z_{k}-\eta\,\tilde g_{k}.
    \]
\EndFor
\State \Return $z_{K}$.
\end{algorithmic}
\end{algorithm}

\section*{Dry Run Example (Univariate Minimisation)}
Consider the simple convex function $f(w)=(w-3)^{2}$, i.e. $L(w)=f(w)$ with no $y$‑variable.  
The gradient is $\nabla f(w)=2(w-3)$.  
We apply the EG scheme treating $y$ as absent, i.e. $z_{k}=w_{k}$ and $F(w)=\nabla f(w)$.  

\begin{center}
\begin{tabular}{c|c|c|c}
Iteration $k$ & $w_{k}$ & $\nabla f(w_{k})$ & $w_{k+1}$\\\hline
0 & $0$ & $2(0-3)=-6$ & $w_{1}=0-\eta[-6-\eta(-6)]=0-0.1(-6-0.1(-6))=0.66$\\
1 & $0.66$ & $2(0.66-3)=-4.68$ & $w_{2}=0.66-0.1[-4.68-0.1(-4.68)]=0.66-0.1(-4.68+0.468)=1.174$\\
2 & $1.174$ & $2(1.174-3)=-3.652$ & $w_{3}=1.174-0.1[-3.652-0.1(-3.652)]=1.174-0.1(-3.652+0.3652)=1.570$\\
\end{tabular}
\end{center}
After three iterations the objective values are
\[
f(0)=9,\qquad f(0.66)\approx5.44,\qquad f(1.174)\approx3.31,\qquad f(1.570)\approx2.06,
\]
showing monotone decrease toward the optimum $w^{\star}=3$.

\newpage
\section*{Assumptions}
\begin{assumption}[Convex–Concave Structure]\label{ass:cc}
For all $x,x'\in\mathcal{X}$ and $y,y'\in\mathcal{Y}$,
\[
L(x',y)-L(x,y)\ge\langle\nabla_{x}L(x,y),x'-x\rangle,\qquad
L(x,y)-L(x,y')\le\langle\nabla_{y}L(x,y),y-y'\rangle .
\]
Thus $L(\cdot,y)$ is convex and $L(x,\cdot)$ is concave.
\end{assumption}

\begin{assumption}[Smoothness]\label{ass:smooth}
There exists $L_{F}>0$ such that the operator $F$ is $L_{F}$‑Lipschitz:
\[
\|F(z)-F(z')\|_{2}\;\le\;L_{F}\|z-z'\|_{2},
\qquad\forall z,z'\in\mathcal{Z}.
\]
Equivalently,
$\|\nabla_{x}L(x,y)-\nabla_{x}L(x',y')\|\le L_{F}\|(x,y)-(x',y')\|$ and the same for $\nabla_{y}$.
\end{assumption}

\begin{assumption}[Existence of a Saddle Point]\label{ass:exist}
There exists $(x^{\star},y^{\star})\in\mathcal{X}\times\mathcal{Y}$ such that
\[
L(x^{\star},y)\le L(x^{\star},y^{\star})\le L(x,y^{\star}),\qquad\forall (x,y)\in\mathcal{X}\times\mathcal{Y}.
\]
Denote $z^{\star}:=(x^{\star},y^{\star})$.
\end{assumption}

\begin{assumption}[Bounded Domain]\label{ass:bounded}
The feasible set $\mathcal{Z}$ is bounded and its diameter is
\[
D:=\max_{z\in\mathcal{Z}}\|z-z^{\star}\|_{2}<\infty .
\]
\end{assumption}

All subsequent results hold under Assumptions \ref{ass:cc}–\ref{ass:bounded}.

\section*{Main Convergence Theorem}
\begin{theorem}[Ergodic $O(1/k)$ rate for the Extra‑Gradient Method]\label{thm:eg}
Let $\{z_{k}\}_{k\ge0}$ be generated by \eqref{eq:mid}–\eqref{eq:update} with a constant stepsize
\[
0<\eta\le\frac{1}{2L_{F}}.
\]
Define the ergodic (averaged) iterate
\[
\bar z_{K}:=\frac{1}{K}\sum_{k=0}^{K-1}z_{k+\frac12}.
\]
Then for any saddle point $z^{\star}$,
\[
\gap(\bar z_{K})\;:=\;\max_{y\in\mathcal{Y}}L(\bar x_{K},y)-\min_{x\in\mathcal{X}}L(x,\bar y_{K})
\;\le\;\frac{D^{2}}{\eta K}.
\]
Consequently $\gap(\bar z_{K})=O(1/K)$.
\end{theorem}

\section*{Theoretical Properties}
\begin{itemize}
\item \textbf{Stability.} The stepsize condition $\eta\le 1/(2L_{F})$ guarantees that the iterates remain in $\mathcal{Z}$ and the distance $\|z_{k}-z^{\star}\|$ is non‑increasing in expectation.
\item \textbf{Hyperparameter Sensitivity.} The bound is linear in $1/\eta$; a larger $\eta$ accelerates convergence up to the limit $1/(2L_{F})$, beyond which the monotonicity argument fails.
\item \textbf{Comparison to Gradient Descent.} For the same class of problems, plain gradient descent yields $O(1/\sqrt{K})$ in the worst case, whereas EG achieves the optimal $O(1/K)$ rate without requiring strong convexity/concavity.
\end{itemize}

\section*{Discussion}
The extra‑gradient method can be interpreted as a predictor–corrector scheme for solving the monotone inclusion $0\in F(z)$. Its $O(1/K)$ ergodic rate matches that of the optimal first‑order methods (e.g., Mirror‑Prox) and is attained under merely Lipschitz continuity of $F$. In practice the two‑step nature often yields superior empirical stability on bilinear games and GAN training.

\newpage
\section*{Preliminaries (Definitions and Lemmas)}
\begin{definition}[Monotone Operator]
An operator $F:\mathcal{Z}\to\R^{n+m}$ is monotone if
\[
\langle F(z)-F(z'),\,z-z'\rangle\ge0,\qquad\forall z,z'\in\mathcal{Z}.
\]
\end{definition}
\begin{lemma}[Monotonicity of $F$]\label{lem:mono}
Under Assumption \ref{ass:cc}, the operator $F$ defined above is monotone.
\end{lemma}
\begin{proof}
For any $z=(x,y),\,z'=(x',y')$,
\[
\begin{aligned}
\langle F(z)-F(z'),z-z'\rangle
&=\langle\nabla_{x}L(x,y)-\nabla_{x}L(x',y'),\,x-x'\rangle\\
&\quad-\langle\nabla_{y}L(x,y)-\nabla_{y}L(x',y'),\,y-y'\rangle.
\end{aligned}
\]
By convexity in $x$ and concavity in $y$ (Assumption \ref{ass:cc}) each term is $\ge0$, hence the sum is $\ge0$.
\end{proof}

\begin{lemma}[Lipschitz Continuity Implies Co‑coercivity of $F$]\label{lem:coercive}
If $F$ is $L_{F}$‑Lipschitz, then for any $z,z'$,
\[
\langle F(z)-F(z'),z-z'\rangle\;\ge\;\frac{1}{L_{F}}\|F(z)-F(z')\|^{2}.
\]
\end{lemma}
\begin{proof}
By Cauchy–Schwarz and the Lipschitz bound,
\[
\|F(z)-F(z')\|^{2}\le L_{F}\langle F(z)-F(z'),z-z'\rangle,
\]
which after rearrangement yields the claim. \qedhere
\end{proof}

\section*{Main Proof (Step‑by‑Step Derivation)}
We prove Theorem \ref{thm:eg}. Throughout the proof we fix an arbitrary saddle point $z^{\star}$.

\paragraph{Step 1 (One‑step expansion).}
From the corrector update \eqref{eq:update} we have
\[
\|z_{k+1}-z^{\star}\|^{2}
= \|z_{k}-\eta F(z_{k+\frac12})-z^{\star}\|^{2}
= \|z_{k}-z^{\star}\|^{2}
-2\eta\langle F(z_{k+\frac12}),\,z_{k}-z^{\star}\rangle
+\eta^{2}\|F(z_{k+\frac12})\|^{2}. \tag{1}
\]
\[
\text{[Step 1]}
\]

\paragraph{Step 2 (Insert predictor relation).}
Subtract and add $\eta F(z_{k})$ inside the inner product:
\[
\langle F(z_{k+\frac12}),z_{k}-z^{\star}\rangle
= \langle F(z_{k}),z_{k}-z^{\star}\rangle
+ \langle F(z_{k+\frac12})-F(z_{k}),z_{k}-z^{\star}\rangle .
\tag{2}
\]
\[
\text{[Step 2]}
\]

\paragraph{Step 3 (Use monotonicity).}
Since $F$ is monotone (Lemma \ref{lem:mono}),
\[
\langle F(z_{k}),z_{k}-z^{\star}\rangle\ge
\langle F(z^{\star}),z_{k}-z^{\star}\rangle=0,
\]
because $F(z^{\star})=0$ at a saddle point. Hence the first term in (2) is non‑negative and can be dropped:
\[
\langle F(z_{k+\frac12}),z_{k}-z^{\star}\rangle
\ge \langle F(z_{k+\frac12})-F(z_{k}),z_{k}-z^{\star}\rangle .
\tag{3}
\]
\[
\text{[Step 3]}
\]

\paragraph{Step 4 (Apply Cauchy–Schwarz and Lipschitz).}
Using Cauchy–Schwarz and the Lipschitz bound of Assumption \ref{ass:smooth},
\[
\langle F(z_{k+\frac12})-F(z_{k}),z_{k}-z^{\star}\rangle
\le \|F(z_{k+\frac12})-F(z_{k})\|\,\|z_{k}-z^{\star}\|
\le L_{F}\|z_{k+\frac12}-z_{k}\|\,\|z_{k}-z^{\star}\|.
\tag{4}
\]
\[
\text{[Step 4]}
\]

\paragraph{Step 5 (Express predictor difference).}
From the predictor step \eqref{eq:mid},
\[
z_{k+\frac12}-z_{k}= -\eta F(z_{k}),
\]
hence $\|z_{k+\frac12}-z_{k}\| = \eta\|F(z_{k})\|$. Substituting into (4) gives
\[
\langle F(z_{k+\frac12})-F(z_{k}),z_{k}-z^{\star}\rangle
\le \eta L_{F}\|F(z_{k})\|\,\|z_{k}-z^{\star}\|.
\tag{5}
\]
\[
\text{[Step 5]}
\]

\paragraph{Step 6 (Bound the squared norm term).}
Using Lemma \ref{lem:coercive},
\[
\|F(z_{k+\frac12})\|^{2}
\le L_{F}\langle F(z_{k+\frac12}),z_{k+\frac12}-z^{\star}\rangle .
\tag{6}
\]
\[
\text{[Step 6]}
\]

\paragraph{Step 7 (Combine (1)–(6)).}
Insert (3)–(5) into (1) and drop the non‑negative term $\langle F(z_{k}),z_{k}-z^{\star}\rangle$:
\[
\begin{aligned}
\|z_{k+1}-z^{\star}\|^{2}
&\le \|z_{k}-z^{\star}\|^{2}
-2\eta\bigl[\langle F(z_{k+\frac12})-F(z_{k}),z_{k}-z^{\star}\rangle\bigr]
+\eta^{2}\|F(z_{k+\frac12})\|^{2}\\
&\le \|z_{k}-z^{\star}\|^{2}
-2\eta\bigl[-\eta L_{F}\|F(z_{k})\|\|z_{k}-z^{\star}\|\bigr]
+\eta^{2}\|F(z_{k+\frac12})\|^{2}\\
&= \|z_{k}-z^{\star}\|^{2}
+2\eta^{2}L_{F}\|F(z_{k})\|\|z_{k}-z^{\star}\|
+\eta^{2}\|F(z_{k+\frac12})\|^{2}.
\end{aligned}
\tag{7}
\]
\[
\text{[Step 7]}
\]

\paragraph{Step 8 (Young’s inequality).}
Apply $ab\le\frac{1}{2}a^{2}+\frac{1}{2}b^{2}$ with $a=\|F(z_{k})\|$ and $b=L_{F}\|z_{k}-z^{\star}\|$:
\[
2\eta^{2}L_{F}\|F(z_{k})\|\|z_{k}-z^{\star}\|
\le \eta^{2}\|F(z_{k})\|^{2}+\eta^{2}L_{F}^{2}\|z_{k}-z^{\star}\|^{2}.
\tag{8}
\]
\[
\text{[Step 8]}
\]

\paragraph{Step 9 (Use Lemma \ref{lem:coercive} on the last term).}
From (6) and the stepsize restriction $\eta\le\frac{1}{2L_{F}}$ we obtain
\[
\eta^{2}\|F(z_{k+\frac12})\|^{2}
\le \eta^{2}L_{F}\langle F(z_{k+\frac12}),z_{k+\frac12}-z^{\star}\rangle
= \eta\langle F(z_{k+\frac12}),\,z_{k+\frac12}-z^{\star}\rangle -\eta^{2}L_{F}^{2}\|z_{k+\frac12}-z^{\star}\|^{2}
\le \eta\langle F(z_{k+\frac12}),\,z_{k+\frac12}-z^{\star}\rangle .
\tag{9}
\]
\[
\text{[Step 9]}
\]

\paragraph{Step 10 (Telescoping inequality).}
Insert (8) and (9) into (7):
\[
\|z_{k+1}-z^{\star}\|^{2}
\le \|z_{k}-z^{\star}\|^{2}
+\eta^{2}\|F(z_{k})\|^{2}
+\eta^{2}L_{F}^{2}\|z_{k}-z^{\star}\|^{2}
+\eta\langle F(z_{k+\frac12}),z_{k+\frac12}-z^{\star}\rangle .
\tag{10}
\]
Rearrange to bring the inner product to the left:
\[
\eta\langle F(z_{k+\frac12}),z_{k+\frac12}-z^{\star}\rangle
\le \|z_{k}-z^{\star}\|^{2}-\|z_{k+1}-z^{\star}\|^{2}
+\eta^{2}\|F(z_{k})\|^{2}
+\eta^{2}L_{F}^{2}\|z_{k}-z^{\star}\|^{2}.
\tag{11}
\]
\[
\text{[Step 10]}
\]

\paragraph{Step 11 (Summation).}
Sum \eqref{11} over $k=0,\dots,K-1$ and use telescoping of the norm terms:
\[
\eta\sum_{k=0}^{K-1}\langle F(z_{k+\frac12}),z_{k+\frac12}-z^{\star}\rangle
\le \|z_{0}-z^{\star}\|^{2}
+\eta^{2}\sum_{k=0}^{K-1}\|F(z_{k})\|^{2}
+\eta^{2}L_{F}^{2}\sum_{k=0}^{K-1}\|z_{k}-z^{\star}\|^{2}.
\tag{12}
\]
\[
\text{[Step 11]}
\]

\paragraph{Step 12 (Bounding the RHS).}
Because $\mathcal{Z}$ has diameter $D$ (Assumption \ref{ass:bounded}), $\|z_{k}-z^{\star}\|\le D$ for all $k$. Moreover, by Lipschitzness,
$\|F(z_{k})\|\le L_{F}D$. Hence
\[
\eta^{2}\sum_{k=0}^{K-1}\|F(z_{k})\|^{2}
\le \eta^{2}K(L_{F}D)^{2},\qquad
\eta^{2}L_{F}^{2}\sum_{k=0}^{K-1}\|z_{k}-z^{\star}\|^{2}
\le \eta^{2}L_{F}^{2}KD^{2}.
\]
Thus the RHS of (12) is bounded by
\[
\|z_{0}-z^{\star}\|^{2}+2\eta^{2}L_{F}^{2}KD^{2}
\le D^{2}+2\eta^{2}L_{F}^{2}KD^{2}.
\tag{13}
\]
\[
\text{[Step 12]}
\]

\paragraph{Step 13 (From monotone operator to primal–dual gap).}
For any $z=(x,y)$ and saddle point $z^{\star}$,
monotonicity of $F$ yields
\[
\langle F(z),z-z^{\star}\rangle
= \langle\nabla_{x}L(x,y),x-x^{\star}\rangle
-\langle\nabla_{y}L(x,y),y-y^{\star}\rangle
\ge L(x,y)-L(x^{\star},y^{\star}) .
\]
Applying this to $z_{k+\frac12}$ and averaging,
\[
\frac{1}{K}\sum_{k=0}^{K-1}\bigl[L(x_{k+\frac12},y_{k+\frac12})-L(x^{\star},y^{\star})\bigr]
\le \frac{1}{\eta K}\sum_{k=0}^{K-1}\langle F(z_{k+\frac12}),z_{k+\frac12}-z^{\star}\rangle .
\tag{14}
\]
\[
\text{[Step 13]}
\]

\paragraph{Step 14 (Combine (12)–(14)).}
Using (12)–(13) in (14) we obtain
\[
\frac{1}{K}\sum_{k=0}^{K-1}\bigl[L(x_{k+\frac12},y_{k+\frac12})-L(x^{\star},y^{\star})\bigr]
\le \frac{D^{2}}{\eta K}+2\eta L_{F}^{2}D^{2}.
\]
Choosing $\eta\le\frac{1}{2L_{F}}$ eliminates the second term:
\[
\frac{1}{K}\sum_{k=0}^{K-1}\bigl[L(x_{k+\frac12},y_{k+\frac12})-L(x^{\star},y^{\star})\bigr]
\le \frac{D^{2}}{\eta K}.
\tag{15}
\]
\[
\text{[Step 14]}
\]

\paragraph{Step 15 (Relation to the primal–dual gap).}
By definition of the saddle‑point,
\[
\gap(\bar z_{K})=
\max_{y}L(\bar x_{K},y)-\min_{x}L(x,\bar y_{K})
\le \frac{1}{K}\sum_{k=0}^{K-1}\bigl[L(x_{k+\frac12},y_{k+\frac12})-L(x^{\star},y^{\star})\bigr],
\]
where $\bar z_{K}=\frac{1}{K}\sum_{k=0}^{K-1}z_{k+\frac12}$. Combining with (15) yields
\[
\gap(\bar z_{K})\le\frac{D^{2}}{\eta K}.
\]
This completes the proof of Theorem \ref{thm:eg}. \qed

\section*{Rate Derivation (With Constants)}
The bound in Theorem \ref{thm:eg} can be written explicitly as
\[
\gap(\bar z_{K})\le \frac{D^{2}}{\eta K},
\qquad\text{with }\eta=\frac{1}{2L_{F}}\;\Longrightarrow\;
\gap(\bar z_{K})\le \frac{2L_{F}D^{2}}{K}.
\]
Hence the constant $C:=2L_{F}D^{2}$ governs the speed of decay.

\section*{Special Cases}
\begin{itemize}
\item \textbf{Bilinear games.} If $L(x,y)=x^{\top}Ay$ with $A$ bounded, then $L_{F}=\|A\|_{2}$ and $D$ is the radius of the feasible sets. The EG method converges linearly to the unique Nash equilibrium when $A$ is positive definite on the primal side and negative definite on the dual side (strong monotonicity).
\item \textbf{Strongly convex–strongly concave.} If $L$ is $\mu$‑strongly convex in $x$ and $\mu$‑strongly concave in $y$, the monotone operator $F$ is $\mu$‑strongly monotone, and the same analysis yields an \emph{exponential} rate $\mathcal{O}((1-\eta\mu)^{K})$.
\item \textbf{Stochastic gradients.} Replacing $F(z_{k})$ by an unbiased estimator with bounded variance adds an extra $\mathcal{O}(1/\sqrt{K})$ term, reproducing the standard $O(1/K)$ ergodic rate for stochastic extra‑gradient methods.
\end{itemize}

\end{document}